\section{Sumcheck algorithms}
\label{app:sumcheck}

This section includes an optimized sumcheck protocol in \cref{fig:sumcheck} and a linear time and space prover algorithm for it in \cref{fig:linear-sc}, both from from~\cite{BagadDaoDombThaler25}. We include these both for self-containedness and because these particular protocols provide a launching off point for some further sumcheck optimizations introduced in our paper.


In the following version of sumcheck the polynomial $s_i$ in each round $i$ is specified by its evaluation on the $d+1$ points $\widehat{U}_d = \set{\infty} \cup \set{0,\dotsc,d}$ (where the point $\infty$ is explained just below). The evaluation $s_i(1)$ can actually be omitted, using the equation $s_i(1) = \sigma_{i-1} - s_i(0)$.

For a degree $d$ univariate polynomial $s(X) = a_d X^d + \dotsm$ we use a variant of Lagrange interpolation where the highest degree coefficient $a_d$ is sent in addition to evaluations of the polynomial $t(X) = s(X) - a_d \cdot \prod_{k=0}^{d-1} (X - x_k)$ over the domain $\set{0,\dotsc,d}$. We note that $t(X)$ is of degree $d-1$, and equal to $s(x)$ over $\set{0,\dotsc,d}$. Defining $s(\infty) = a_d$, we recover $s$ via evaluations of
\[
s(X) = a_d \cdot \prod_{k=1}^d (X - x_k) + \sum_{k=1}^d t(x_k) \cdot \mathcal{L}_{\set{0,\dotsc,d},k}(X)
\]
over $\widehat{U}_d$, where $\mathcal{L}_{\set{0,\dotsc,d},k}$ is the $k$th Lagrange basis polynomial over the interpolation set $\set{0,\dotsc,d}$.

If $s(X) = \prod_{i=1}^d p_i(X)$ is a product of linear polynomials, then
\[
s(\infty) = \prod_{i=1}^d p_i(\infty) = \prod_{i=1}^d (p_i(1) - p_i(0)) \quad.
\]
More generally, when $s(X) = G(p_1(X), \ldots, p_d(X))$ for some polynomial $G$ and linear polynomials $p_1, \ldots, p_d$, we can calculate $s(\infty)$ by evaluating only the terms of $G$ of highest total degree. Then, like $s(1)$, the prover can generally omit sending the evaluation $s(\infty)$, instead computing it using the equation above.

\begin{figure}[h!]
\centering
\caption{The sumcheck protocol}
\label{fig:sumcheck}
\vspace{-1em}
\begin{mdframed}[linewidth=1pt,roundcorner=8pt,backgroundcolor=gray!2,innerleftmargin=10pt,innerrightmargin=10pt]
\textbf{Setup}: Assume that the verifier has oracle access to the polynomial $g$, and knows the value $\sigma_0$ claimed to equal
\[
\sum_{x\in\{0,1\}^v} g(x) = \sigma_0.
\]

\textbf{For each round} $i = 1, \ldots, v$:
\begin{enumerate}
    \item The prover $\prover$ sends the univariate polynomial $s_i(X)$ claimed to equal
    \[
    \sum_{(x_{i+1},\ldots,x_v)\in\{0,1\}^{v-i}} g(\alpha_1,\ldots,\alpha_{i-1}, X, x_{i+1},\ldots,x_v).
    \]
    $\prover$ does so by sending the evaluations $\{ s_i(u) : u \in \widehat{U}_d \}$ to the verifier $\verifier$.
    \item The verifier $\verifier$ sends a random challenge $\alpha_i \leftarrow \mathbb{F}$ to $\prover$.
    \item $\verifier$ derives $s_i(1) := \sigma_{i-1} - s_i(0)$, then sets $\sigma_i := s_i(\alpha_i)$.
\end{enumerate}

After $v$ rounds, we reduce to the claim that
\[
g(\alpha_1, \ldots, \alpha_v) = \sigma_v.
\]
$V$ checks this claim by making a single oracle query to $g$.
\end{mdframed}
\end{figure}

A linear time and space prover for this sumcheck version may be implemented as follows. The prover maintains $d$ arrays $P_1, \dotsc, P_d$, which initially store all evaluations of $p_1, \dotsc, p_d$ over $\bits^v$; we have $P_k[x] = p_k(x)$ for all $k \in [d]$ and $x \in \bits^v$. The prover halves the size of the arrays after each round $i$, via “binding” the arrays $[P_k]_{k\in [d]}$ to the challenge $\alpha_i$ received from the verifier.

\begin{figure}[h!]
\centering
\caption{A linear time and space prover implementation of sumcheck}
\label{fig:linear-sc}
\vspace{-1em}
\begin{mdframed}[linewidth=1pt,roundcorner=8pt,backgroundcolor=gray!2,innerleftmargin=10pt,innerrightmargin=10pt]

\textbf{Setup}: Length-$2^{v}$ arrays $P^{(0)}_1,\ldots,P^{(0)}_d$ such that
\[
P^{(0)}_k[x] = p_k(x) \ \text{for all } k = 1,\ldots,d \ \text{and } x \in \{0,1\}^{v}.
\]

\textbf{For each round} $i = 1,\ldots,v$:
\begin{enumerate}
  \item For $u \in U_{d}^{b}$, compute $s_i(u)$ using the formula
  \[
  s_i(u) := \sum_{x' \in \{0,1\}^{v-i}} \prod_{k=1}^{d}
  \Big( \big(P^{(i)}_k[1,x'] - P^{(i)}_k[0,x']\big)\cdot u + P^{(i)}_k[0,x'] \Big).
  \]
  \item Send $\{\, s_i(u) : u \in U_{d}^{b} \,\}$ to the verifier.
  \item Receive challenge $\alpha_i \in F$ from the verifier.
  \item For $k = 1,\ldots,d$ and $x' \in \{0,1\}^{v-i}$, update arrays:
  \[
  P^{(i+1)}_k[x'] := \big(P^{(i)}_k[1,x'] - P^{(i)}_k[0,x']\big)\cdot \alpha_i + P^{(i)}_k[0,x'].
  \]
\end{enumerate}

\texttt{// Invariant: }$P^{(i)}_k[x'] = p_k(r[1:i], x')$ for all $k \in [d]$, $i \in [0,v]$, $x' \in \{0,1\}^{v-i}$.

\end{mdframed}
\end{figure}